\documentclass[12pt,a4paper]{article}

\usepackage[utf8]{inputenc}
\usepackage[T1]{fontenc}
\usepackage[margin=2.5cm]{geometry}
\usepackage{hyperref}
\usepackage{microtype}
\usepackage{parskip}
\usepackage{lmodern}

\hypersetup{colorlinks=true, urlcolor=blue!60!black}

\pagestyle{empty}

\begin{document}

\begin{flushright}
Kundai Farai Sachikonye \\
Auf der Lichtung 19, 82178 Puchheim, Germany \\
\href{https://github.com/fullscreen-triangle}{github.com/fullscreen-triangle} \\
\texttt{kundai.sachikonye@bitspark.com} \\
(+49) 1626386344 \\[6pt]
20 May 2026
\end{flushright}

\vspace{1em}

Frau Lisa Merkel \\
Pathologisches Institut --- Zentrum für Personalisierte Medizin (ZPM) \\
Universitätsklinikum Heidelberg \\
Im Neuenheimer Feld 224, 69120 Heidelberg \\
Job-ID: V000015540

\vspace{1.5em}

\textbf{Re: Data Manager*in --- Onkologie, ZPM Heidelberg (V000015540)}

\vspace{1em}

Dear Frau Merkel,

I am applying for the Data Manager position (V000015540) at the ZPM Heidelberg.
My interest in this role is specific: the ZPM operates at the interface between
clinical oncology and translational research, and that is precisely the context
in which the data management problems are hardest.
In purely clinical settings, data flows are standardised.
In purely research settings, data structures are investigator-defined.
At the interface, neither set of conventions applies cleanly, and managing
data quality requires understanding what the data will be used for --- not just
that it has been entered correctly.

\textbf{Oncology data infrastructure at clinical scale.}
Over the past four years I have built the analytical layer that sits downstream
of clinical data collection in oncology.
The \href{https://syndrome-seven.vercel.app/tools}{Syndrome} framework processes
patient and tumour characterisation data to produce neoantigen predictions,
MHC binding assessments, immune escape scores, and a structured disease state
representation $\mathbf{D}\in[0,1]^8$ with dimensions including immune activation,
mucosal barrier integrity, metabolic, and circadian components.
The \href{https://github.com/fullscreen-triangle/gospel}{Gospel} pipeline
handles genomic variant calling, somatic mutation analysis, RNA-seq quantification,
and multi-omics integration from clinical cohort data, validated against
gnomAD, 1000~Genomes, and UK~Biobank at $10^6$ variants per second.
The \href{https://github.com/fullscreen-triangle/lavoisier}{Lavoisier} pipeline
processes mass spectrometry data ($>100$\,GB mzML) from clinical samples with
98.9\,\% NIST classification accuracy.
Building these systems required understanding, at each step, how clinical
data is generated, what the quality failure modes are, and what the downstream
analysis requires from the upstream documentation.
That is the same understanding the ZPM data manager role requires in the
other direction: knowing what the research and certification use cases need
in order to maintain the right data quality upstream.

\textbf{Clinical data experience.}
At Universitätsklinikum Hamburg-Eppendorf (2017--2018) I worked directly
in the Glycomics and Proteomics laboratory, processing clinical samples,
operating automated LC-MS/MS and MALDI-TOF pipelines, and managing data
annotation and quality control for clinical cohort studies.
At TUM / Bayerisches Forschungsinstitut (2019--2021) I worked on
multi-omics data integration across federated clinical sites, maintaining
distributed data processing pipelines and managing data quality across
heterogeneous institutional data structures.
I am familiar with the practical constraints of clinical data: incomplete
records, inconsistent coding, retrospective vs.\ prospective collection
differences, and the documentation requirements for certification and
reporting that differ from what researchers want.

\textbf{eCRF development and EDC systems.}
I do not have prior experience with Onkostar specifically, and I want to
be direct about that.
What I bring instead is deep familiarity with the data structures
that tumour documentation systems capture and the analysis pipelines
that consume them.
eCRF design is a data modelling problem: the form must capture exactly
what the downstream analysis requires, in a structure that clinicians
can complete reliably under time pressure.
I have made those modelling decisions from the analysis side many times,
and I expect the learning curve on Onkostar's specific interface to be short.

\textbf{Scientific publications and reporting.}
The position includes contributing to scientific publications and
certification reporting.
I have written formal manuscripts with full mathematical apparatus ---
theorem statements, validation experiments, quantitative performance claims
--- and I am comfortable with the standards of precision clinical research
publications require.
For the ZPM context, this means I can contribute to translational
manuscripts that span clinical data, genomics, and computational analysis,
not just to methods sections.

\textbf{Interdisciplinary coordination.}
The ZPM brings together pathologists, oncologists, computational researchers,
and IT infrastructure teams.
The data manager sits at the centre of that coordination.
My work has required continuous communication across exactly this range
of expertise: explaining computational methods to clinicians, explaining
clinical constraints to software engineers, and translating between the
two in project planning and documentation.
User training and support --- another component of this role --- are familiar
from deploying tools that wet-lab and clinical collaborators use without
local computational infrastructure.

\textbf{Background.}
I hold an MSc in Computational Biology (Universität Tübingen), an MSc in
Pharmaceutical Biotechnology (HAW Hamburg), and a BSc in Biochemistry and
Cell Biology (Jacobs University Bremen).
I conducted doctoral research in Computational Lipidomics at TUM (2019--2021).
I have lived in Germany continuously since 2011.
English is native; German is conversational.
I am available immediately.

\vspace{1.5em}
Yours sincerely,

\vspace{2.5em}
\textbf{Kundai Farai Sachikonye}

\end{document}
